\section{从边带和式算出阿秒时间密度}
\label{chap:feqo-attosecond}

光学周期 \(T_\omega=2\pi/\omega\) 内的归一化时间密度定义为 \(I(\tau;L)=|\psi(\tau;L)|^2/T_\omega\)，其中 \(\psi\) 取式~\eqref{eq:feqo-drift-envelope} 且 \(\sum|c_\ell|^2=1\)。这里把绝对到达时刻按模 \(T_\omega\) 折叠；它是一个周期内的相位或到达时间分布，不是宣称无限长的电子波列在整个实时间轴上的概率积分为一。若要预测有限脉冲的绝对到达分布，必须保留前节的连续能量包络。展开模平方，
\begin{equation}
I(\tau;L)=\frac1{T_\omega}\sum_{\ell,\ell'}c_\ell c_{\ell'}^*
e^{-\ii(\ell-\ell')\omega\tau}
e^{\ii(\Phi_\ell-\Phi_{\ell'})}.
\label{eq:feqo-temporal-density}
\end{equation}
周期积分中的 \(e^{-\ii(\ell-\ell')\omega\tau}\) 正交，故 \(\int_0^{T_\omega}I\dd\tau=\sum|c_\ell|^2=1\)。这项检查说明时间峰只是概率在一个周期内重新分配，没有创造电子。

要把能谱与时间图形的差别定量说清，按整数 \(m\) 收集式~\eqref{eq:feqo-temporal-density} 中频率为 \(m\omega\) 的项，定义
\begin{equation}
b_m(L):=\sum_{\ell\in\mathbb Z}c_\ell c_{\ell-m}^*
e^{\ii[\Phi_\ell(L)-\Phi_{\ell-m}(L)]},
\qquad
T_\omega I(\tau;L)=\sum_{m\in\mathbb Z}b_m(L)e^{-\ii m\omega\tau}.
\label{eq:feqo-bunching-harmonics}
\end{equation}
\(b_0=1\)、\(b_{-m}=b_m^*\)，故密度为实数且已归一。由柯西--施瓦茨不等式，\(|b_m|\leq(\sum|c_\ell|^2)^{1/2}(\sum|c_{\ell-m}|^2)^{1/2}=1\)。它是相邻 \(m\) 格的\emph{相干振幅}之和；只知道 \(P_\ell=|c_\ell|^2\) 不能求它。若任一相关边带相位被随机化，各项在重复实验中互相抵消，\(|b_m|\) 就变小。

对理想 PINEM 相位调制，交互区出口处的周期波函数是模长为一的 \(u(\theta)=e^{\beta e^{-\ii\theta}-\beta^*e^{\ii\theta}}\)。把 \(|u|^2=1\) 作 Fourier 展开，立即得到 \(b_m(0)=0\)（\(m\ne0\)）。因此边带在出口已经存在，时间密度却仍平坦；随后出现的时间尖峰不是光场直接把电子“切成脉冲”，而是自由色散使原先相消的跨边带项取得不同相位。

为看见峰怎样长出来，暂时只保留 \(\ell=-1,0,1\)，设 \(c_0=J_0(2b)\)、\(c_{+1}=J_1(2b)e^{\ii\phi}\)、\(c_{-1}=-J_1(2b)e^{-\ii\phi}\)，其中 \(b=|\beta|\)。在二次色散下令 \(\Phi_{\pm1}=\Phi_2/2\)、\(\Phi_0=0\)，则
\[
\psi_3=J_0(2b)-2\ii J_1(2b)e^{\ii\Phi_2/2}
\sin(\omega\tau-\phi),
\]
从而
\begin{align*}
T_\omega I_3={}&J_0^2(2b)+4J_1^2(2b)\sin^2(\omega\tau-\phi)\\
&+4J_0(2b)J_1(2b)\sin(\Phi_2/2)\sin(\omega\tau-\phi).
\end{align*}
若 \(\Phi_2=0\)，最后一项消失，弱耦合时密度变化从 \(b^2\) 才开始；传播到 \(|\Phi_2|\simeq\pi\) 时，一阶谐波的幅度达到本三边带模型的最大值。三边带截断的积分是 \(J_0^2+2J_1^2\)，与一的差来自舍弃的 \(|\ell|\geq2\) 人口，量级为 \(O(b^4)\)。强耦合时须恢复完整和式。

写 \(\Phi_\ell\simeq\tfrac12\ell^2\Phi_2\)，并由式~\eqref{eq:feqo-sideband-phase-expansion} 得 \(\Phi_2=-\hbar\omega^2L/(\gamma_0^3mv_0^3)\)。当 \(|\Phi_2|=2\pi\) 时，因 \(\ell^2\equiv\ell\pmod2\)，二次相位 \(e^{\ii\pi\ell^2}=e^{\ii\pi\ell}\) 只把时间图形平移半个周期。因此时间密度的 Talbot 复现尺度为
\begin{equation}
L_T=\frac{2\pi\gamma_0^3mv_0^3}{\hbar\omega^2}.
\label{eq:feqo-talbot-length}
\end{equation}
完整波函数不带时间平移地复现要到 \(2L_T\)。这项结论仍受三阶余项控制；若 \(\max|R_{3,\ell}^{(p)}|\gtrsim1\)，理想复现会被破坏。

初始能散、相位抖动与距离抖动要分别进入。若电子中心能量有宽度 \(\sigma_E\) 的窄高斯分布，不同电子的群延迟差一阶为 \(\delta\tau=-L\delta E/(\gamma_0^3mv_0^3)\)；能量较高的电子先到达。对到达时间作经典平均后，第 \(m\) 个时间谐波乘
\[
\exp\!\left[-\frac12\left(\frac{m\omega L\sigma_E}{\gamma_0^3mv_0^3}\right)^2\right].
\]
若每次制备的光学相位独立地以方差 \(\sigma_\phi^2\) 抖动，另一因子为 \(e^{-m^2\sigma_\phi^2/2}\)。距离分布则需对 \(\Phi_\ell(L)\) 再平均。这三种操作物理来源不同，不能用一个“模糊宽度”代替。

\begin{exercise}
对三边带公式积分一个周期，并计算 \(\Phi_2=0\) 与 \(\Phi_2=\pi\) 时的一阶时间谐波幅度。
\end{exercise}
